\usepackage{amsmath, amssymb}
\usepackage{cancel}
\usepackage{siunitx}
\usepackage[T1]{fontenc}
\usepackage{bm}
\usepackage{tikz}
\usetikzlibrary{calc, angles, quotes}
\usetikzlibrary{decorations.pathreplacing}
\usepackage{parskip}
\usepackage[french]{babel}
\usepackage{csquotes}
\MakeOuterQuote{"}
\usepackage{tcolorbox}
\usepackage{pgfplots}

% Pour avoir sum() dans les graphes TikZ
\usepackage{pgfplots}
\pgfplotsset{compat=1.18}

\date{\center \footnotesize Version du \today}

\usepackage{titling}
\setlength{\droptitle}{-2em}      % remonte tout le bloc titre

\usepackage{titlesec}

\titleformat{\paragraph} % On veut que paragraph se comporte comme un titre avec un saut de ligne
  {\normalfont\normalsize\bfseries}
  {\theparagraph}
  {1em}
  {}
  
\usepackage{titlesec}

\tcbuselibrary{theorems,skins,breakable}

\newtcbtheorem[no counter] %auto counter, number within = section]
{definition}{Définition}{
  breakable,
  fonttitle = \bfseries,
  colframe = yellow!25!gray,
  colback = yellow!2,
  boxrule=0.8pt
}{def}

\usepackage[
  a4paper,
  inner=2cm,
  outer=1.5cm,
  top=2cm,
  bottom=2cm
]{geometry}

% Pour l'espacement entre les énumérations qui sont par défaut tassées
\usepackage[inline]{enumitem}
\setlist[enumerate, 1]{topsep=0.75em, itemsep=0.75em}
\setlist[enumerate, 2]{topsep=0.25em, itemsep=0.25em}

% Symbole mathématique des multiensembles {{2, 3, 3, 5, 6}}
\newcommand{\mset}[1]{\{\!\!\{#1\}\!\!\}}
